\section{Tutorial de implementação}

\subsection{Instalação}

Em Debian ou Ubuntu, Open MPI e a interface Python podem ser instalados com:

\begin{lstlisting}[style=terminal]
sudo apt update
sudo apt install openmpi-bin python3-mpi4py
\end{lstlisting}

O programa usa \texttt{MPI.COMM\_WORLD}, que reúne todos os processos
iniciados pelo \texttt{mpiexec}. Cada processo consulta seu identificador
(\textit{rank}) e o número total de processos. A documentação oficial do
\texttt{mpi4py} apresenta esse modelo de execução e as operações coletivas
\cite{mpi4py}.

\subsection{Código}

A função \texttt{calcular\_pi} sorteia os pontos e retorna apenas a quantidade
que pertence ao círculo. Na parte paralela, cada processo recebe
aproximadamente $N/p$ pontos. A operação \texttt{reduce}, com
\texttt{MPI.SUM}, soma os acertos e entrega o total ao processo zero.

\lstinputlisting[
    style=python,
    caption={Implementação serial e MPI em um único arquivo.},
    label={lst:codigo}
]{pi.py}

\subsection{Execução}

O programa é executado com quatro processos, como solicitado:

\begin{lstlisting}[style=terminal]
mpiexec -n 4 python3 pi.py
\end{lstlisting}

O \texttt{mpiexec} é o comando padrão para iniciar os processos MPI
\cite{openmpi}. Caso o Open MPI emita avisos sobre interfaces TCP em uma
execução inteiramente local, pode-se selecionar apenas os transportes locais:

\begin{lstlisting}[style=terminal]
mpiexec --mca btl self,sm -n 4 python3 pi.py
\end{lstlisting}
